\subsection{ASSOCIATIVE LAWS}

DEFINITION 5. A law of composition \((x,y)\mapsto x\top y\) on a set \(\mathbf{E}\) is called associative if, for all elements \(x,y,z\) in \(\mathbf{E}\) ,

\[
(x \top y) \top z = x \top (y \top z).
\]

A magma whose law is associative is called an associative magma.

The opposite law of an associative law is associative.

Examples. (1) Addition and multiplication of natural numbers are associative laws of composition on \(\mathbf{N}\) (Set Theory, III, § 3, no. 3, Corollary to Proposition 5)

\begin{enumerate}
\def\labelenumi{(\arabic{enumi})}
\setcounter{enumi}{1}
\tightlist
\item
  The laws cited in Examples (1), (3) and (4) of no. 1 are associative.
\end{enumerate}

THEOREM 1 (Associativity theorem). Let E be an associative magma whose law is denoted by \(\top\) . Let A be a totally ordered non-empty finite set, which is the union of an ordered sequence of non-empty subsets \((\mathbf{B}_l)_{i \in I}\) such that the relations \(\alpha \in \mathbf{B}_i\) , \(\beta \in \mathbf{B}_j\) , \(i < j\) imply \(\alpha < \beta\) ; let \((x_\alpha)_{\alpha \in A}\) be an ordered sequence of elements in E with A as indexing set. Then

\[
\bigcap_ {\alpha \in A} x _ {\alpha} = \bigcap_ {i \in I} \left(\bigcap_ {\alpha \in B _ {i}} x _ {\alpha}\right)\tag{3}
\]

† The use of this terminology and the notation \(\prod_{\alpha \in \mathbf{A}} x_{\alpha}\) must be avoided if there is any risk of confusion with the product of the sets \(x_{\alpha}\) defined in the theory of sets (Set Theory, II, § 5, no. 3). However, if the \(x_{\alpha}\) are cardinals and addition (resp. multiplication) is the cardinal sum (resp. the cardinal product), the cardinal denoted by \(\sum_{\alpha \in \mathbf{A}} x_{\alpha}\) (resp. \(\prod_{\alpha \in \mathbf{A}} x_{\alpha}\) ) in the above notation is the cardinal sum (resp. cardinal product) of the family \((x_{\alpha})_{\alpha \in \mathbf{A}}\) (Set Theory, III, § 3, no. 3).

We prove the theorem by induction on the cardinal n of A. Let p be the cardinal of I and h its least element; let \(J = I - \{h\}\) . If n = 1, as the \(B_{i}\) are non-empty, of necessity p = 1 and the theorem is obvious. Otherwise, assuming the theorem holds for an indexing set with at most n - 1 elements, we distinguish two cases:

\begin{enumerate}
\def\labelenumi{(\alph{enumi})}
\tightlist
\item
  \(B_{h}\) has a single element \(\beta\) . Let \(C = \bigcup_{i \in J} B_{i}\) . The left-hand side of (3) is equal, by definition, to \(x_{\beta} \top \left( \bigcap_{\alpha \in C} x_{\alpha} \right)\) ; the right-hand side is equal, by definition, to
\end{enumerate}

\[
x _ {\beta} \top \left(\underset {i \in J} {\mathsf {T}} \left(\underset {\alpha \in B _ {i}} {\mathsf {T}} x _ {\alpha}\right)\right);
\]

the result follows from the fact that the theorem is assumed true for C and \((\mathbf{B}_{i})_{i\in J}\) .

\begin{enumerate}
\def\labelenumi{(\alph{enumi})}
\setcounter{enumi}{1}
\tightlist
\item
  Otherwise, let \(\beta\) be the least element of A (and hence of \(\mathbf{B}_h\) ); let \(\mathbf{A}' = \mathbf{A} - \{\beta\}\) and let \(\mathbf{B}_i' = \mathbf{A}' \cap \mathbf{B}_i\) for \(i \in \mathbf{I}\) ; then \(\mathbf{B}_i' = \mathbf{B}_i\) for \(i \in \mathbf{J}\) . The set \(\mathbf{A}'\) has \(n - 1\) elements and the conditions of the theorem hold for \(\mathbf{A}'\) and its subsets \(\mathbf{B}_i'\) ; hence by hypothesis:
\end{enumerate}

\[
\underset {\alpha \in A ^ {\prime}} {T} x _ {\alpha} = \left(\underset {\alpha \in B ^ {\prime} h} {T} x _ {\alpha}\right) \top \left(\underset {i \in J} {T} \left(\underset {\alpha \in B _ {i}} {T} x _ {\alpha}\right)\right).
\]

Forming the composition of \(x_{\beta}\) with each side, we have on the left-hand side side, by definition, \(\bigcap_{\alpha \in A} x_{\alpha}\) and on the right-hand side, using the associativity,

\[
\left(x _ {\beta} \top \left(\underset {\alpha \in B ^ {\prime} h} {\top} x _ {\alpha}\right)\right) \top \left(\underset {i \in J} {\top} \left(\underset {\alpha \in B _ {i}} {\top} x _ {\alpha}\right)\right)
\]

which is equal, by Definition 3, to the right-hand side of formula (3).

For an associative law denoted by \(\top\) the composition \(\bigcap_{p\leqslant i\leqslant q}x_i\) of a sequence \((x_i)_{i\in [p,q]}\) is also denoted (since no confusion can arise) by

\[
x _ {p} \top \dots \top x _ {q}.
\]

A particular case of Theorem 1 is the formula

\[
x _ {0} \top x _ {1} \top \dots \top x _ {n} = (x _ {0} \top x _ {1} \top \dots \top x _ {n - 1}) \top x _ {n}.
\]

Consider an ordered sequence of n terms all of whose terms are equal to the same element \(x \in E\) . The composition of this sequence is denoted by \(\top^{n} x\) for a law denoted by \(\top, \bot^{n} x\) for a law denoted by \(\bot\) . For a law written multiplicatively the composition is denoted by \(x^{n}\) and called the n-th power of x. For a law written additively the composition is usually denoted by nx. The associativity theorem applied to an ordered sequence all of whose terms are equal gives the equation

\[
\begin{array}{c} n _ {1} + n _ {2} + \dots + n _ {p} \\ \top \end{array} x = \binom {n _ {1}} {\top} x) \top \binom {n _ {2}} {\top} x) \top \dots \top \binom {n _ {p}} {\top} x).
\]

In particular, if \(p = 2\) ,

\[
{ } ^ { m + n } \top x = \left( \begin{array} { c } m \\ \top x \end{array} \right) \top \left( \begin{array} { c } n \\ \top x \end{array} \right)\tag{4}
\]

and if \(n_1 = n_2 = \cdots = n_p = m,\)

\[
\frac {p m}{T} x = \frac {p}{T} \left(\frac {m}{T} x\right).\tag{5}
\]

If X is a subset of E, we sometimes denote, in conformity with the above notation, by \(\overset{p}{\top}X\) the set \(X_{1}\top X_{2}\top\cdots\top X_{p}\) , where

\[
\mathbf {X} _ {1} = \mathbf {X} _ {2} = \dots = \mathbf {X} _ {p} = \mathbf {X};
\]

it is thus the set of all compositions \(x_{1} \top x_{2} \top \cdots \top x_{p}\) with \(x_{1} \in X, x_{2} \in X, \ldots, x_{p} \in X\) .

It is important not to confuse this set with the set of \(\top x\) , where x runs through X.

